\chapter{Related Work}
\label{ch:literature}

\section{Real-Time Object Detection}
\label{sec:lit_objdet}

The detector used in this thesis is YOLO~\cite{redmon2016yolo} as packaged in the Ultralytics framework~\cite{ultralytics2024}. The nano variant is chosen for latency: the cascade runs RGB and IR detectors back-to-back plus a trust classifier and a per-frame distilled MLP verifier, leaving little room for a heavier backbone. Two-stage architectures such as Faster R-CNN~\cite{ren2015fasterrcnn} and Cascade R-CNN~\cite{cai2018cascade} reach higher accuracy at a latency cost this budget does not allow. The transformer-based detection family --- DETR~\cite{carion2020detr} and its real-time variant RT-DETR~\cite{zhao2024rtdetr} --- is a plausible drop-in for future work; it was not adopted here because the cascade's downstream stages depend on the YOLO confidence-score distribution against which the trust classifier is calibrated, making backbone replacement a long-horizon decision rather than a free tuning knob.

A frequently under-reported hyperparameter is the inference resolution \texttt{imgsz}. YOLO letterboxes the input to a square canvas; objects smaller than the effective receptive field are unresolvable regardless of model quality. On Svanstr\"om (native $640\times512$) the effect is stark: at \texttt{imgsz=1280} the baseline reaches drone $R=0.961$, whereas at \texttt{imgsz=640} small drones fall below the resolvable floor --- the \texttt{retrained\_v2} variant collapses to $R=0.072$ at that size (\textsc{Evidence\_Ledger}~\S3.1, \S10). The general issue of small-object recall in modern detectors is treated by \citeauthor{rozantsev2017detecting}~\cite{rozantsev2017detecting} for the related flying-object case.
% [source: ledger=selcom-imgsz-win; eval=rgb_svan_baseline; run=eval/diagnose_failures.py; cache=eval/results/_failure_diagnosis/svanstrom_1280_by_category.csv; config=svan_iop_1280]

\section{Drone Detection and the Confuser Problem}
\label{sec:lit_drone}

At the ranges typical of aerial surveillance, drones, birds, fixed-wing aircraft, and helicopters all present as small dark silhouettes against bright sky with similar aspect ratios and little fine texture. A YOLO detector trained with bird negatives still fires on 94.4\% of bird-only Svanstr\"om frames (\textsc{Evidence\_Ledger}~\S3.1). This is the \emph{confuser problem}: training-time hard-negative mining suppresses helicopters and partially airplanes but cannot drive bird fire rate down without sacrificing recall on small-drone targets that share the same silhouette --- the one stance that suppresses birds (\texttt{retrained\_v2}) collapses Svanstr\"om drone recall to $R=0.306$.
% [source: ledger=retrainedv2-recall-collapse; eval=rgb_svan_baseline; run=eval/diagnose_failures.py; cache=eval/results/_failure_diagnosis/svanstrom_1280_by_category.csv; config=svan_iop_1280]

C-UAS sensing surveys~\cite{shi2018counteruas,taha2019drone,samaras2019deep} converge on the same conclusion: visual detection is operationally attractive but limited by confuser discrimination at long range. Early deep-learning approaches~\cite{aker2017using,rozantsev2017detecting} treated drone detection as a single-class small-object problem; subsequent work explicitly added confuser classes. \citet{schumann2017deep} report cross-domain bird-vs-UAV classification with deep CNNs at small scale; the Drone-vs-Bird challenge~\cite{coluccia2021dronevsbird} consolidated the bird-confuser benchmark and validated that hard-negative mining is necessary but not sufficient.

Two public benchmarks anchor the empirical work. The Svanstr\"om RGB+IR dataset~\cite{svanstrom2021real,svanstrom2022dronedataset} provides per-frame box annotations with three explicit confuser classes (birds, airplanes, helicopters); its native $640\times512$ resolution forces the small-drone regime, making it the discriminating benchmark. The Anti-UAV RGBT benchmark~\cite{jiang2021antiuav,zhao2023antiuav} provides large-scale paired RGB+IR sequences but is saturated: baseline and \texttt{retrained\_v2} RGB both reach $P=0.9922, R=0.9950, F1=0.9936$ at \texttt{imgsz=1280} ($3{,}178$~TP, $25$~FP, $16$~FN), and IR-only reaches $F1=0.9654$ (\textsc{Evidence\_Ledger}~\S3.2), because Anti-UAV RGB frames appear in both variants' training corpus ($59{,}413$ \texttt{anti}-prefixed frames, Table~\ref{tab:ds_rgb_components}). Anti-UAV is therefore an in-distribution sanity floor, not a discriminating benchmark. \S\ref{sec:comparison} compares the production stack to published numbers on both datasets.
% [source: ledger=antiuav-saturated; eval=rgb_antiuav_baseline; run=eval/ablate.py; cache=eval/results/_ablation/2026-05-18T22-48-19/master.csv; config=antiuav_iou_1280_s5]

The cascade is the response to the residual confuser fire rate that training-time mining cannot eliminate; the trust classifier and per-frame verifier catch what the detector cannot.

\section{Thermal Infrared Imaging}
\label{sec:lit_ir}

Thermal infrared cameras image surface emission rather than reflected visible light, with three properties relevant here. First, drone motors, ESCs, and batteries produce localised hot spots against sky, whereas birds tend to maintain more uniform body temperatures --- a discriminative cue absent in the visible spectrum and a likely reason the IR detector hallucinates less on bird confusers (\S\ref{sec:grayscale}). Second, IR is largely illumination-invariant, aiding night and low-visibility operation. Third, single-channel IR imagery shares low-level statistical properties with grayscale RGB: both render small drones as silhouettes against sky. This shared structure underwrites the cross-modal transfer result in \S\ref{sec:grayscale}, where the IR detector reaches drone $F1=0.636$ on grayscale-converted RGB video and ties the RGB baseline on the hardest bird-cluttered clip. The transfer was emergent, not designed; the grayscale-RGB mode exists as a fallback for when thermal hardware is unavailable.
% [source: ledger=ir-grayscale-fallback; eval=vid_drone_ir_gray; run=eval/eval_video_tests.py; config=video_drone_iop]

IR imaging also has limitations: lower spatial resolution, thermal contrast collapse when ambient approaches drone surface temperature, and smaller labelled datasets complicating training and OOD evaluation. The Roboflow OOD audit (\textsc{Ledger}~\S8.3) makes IR brittleness explicit: $R=0.519$ on \texttt{ir\_mixed\_cbam} and $R=0.264$ on \texttt{ir\_drone\_night}, against $R\geq 0.94$ on Anti-UAV and Svanstr\"om.
% [source: ledger=ir-ood-recall-ceiling; eval=ir_ood_cbam_raw; run=(EVIDENCE_LEDGER sweep); config=roboflow_ir_drone_640]

\section{Multi-Modal Fusion and Hard-Negative Mining}
\label{sec:lit_fusion}

Multi-modal C-UAS systems combine RGB and thermal streams through early, intermediate, or decision-level fusion~\cite{ramachandram2017fusion}. Intermediate-level fusion --- the canonical example being multispectral pedestrian detection~\cite{wagner2016multispectral} --- typically requires paired-modality training data. This thesis instead adopts decision-level fusion with a per-frame learned XGBoost trust classifier~\cite{chen2016xgboost,friedman2001greedy}: the two detectors hallucinate on different scenes, so a per-frame decision picks the modality whose failure profile is least active. Hard-negative mining (OHEM~\cite{shrivastava2016ohem}, Focal Loss~\cite{lin2017focal}, explicit confuser augmentation) suppresses helicopters and airplanes well ($66.2\% \to 41.9\%$ and $74.6\% \to 64.7\%$ for \texttt{hardneg\_v3more}) but leaves bird fire rate essentially unchanged at $94.2\%$, and the \texttt{retrained\_v2} stance that does suppress birds (to $3.4\%$) collapses drone recall to $R=0.306$. Birds and small drones are not separable at the detection stage without an unacceptable recall cost; the downstream cascade is the response to that case.
% [source: ledger=retrainedv2-recall-collapse; eval=rgb_svan_hardneg; run=eval/diagnose_failures_all.py; cache=eval/results/_failure_diagnosis/svanstrom_1280_by_category.csv; config=svan_iop_1280]

\section{Detection Scoring}
\label{sec:lit_calibration}

The choice of matching rule is consequential and under-discussed. Standard IoU under-counts predictions that are correct in centre but imprecise in scale; for small drones whose GT boxes exceed the visible target (Svanstr\"om), this is routine. This thesis uses Intersection over Prediction (IoP) at threshold $0.5$ for Svanstr\"om RGB scoring. A second axis --- dual-modality versus trust-aware F1 aggregation --- swings reported F1 by 28~pp on Svanstr\"om for identical underlying detections (\S\ref{sec:scoring_audit}; \textsc{Ledger}~\S1). Multi-modal drone-detection numbers are not comparable across studies without explicit scoring disclosure.

\section{Cross-Modal Transfer and Cascaded Rejection}
\label{sec:lit_xmodal_cascade}

Two literatures bear directly on the architectural pieces that do not fit neatly under object detection, thermal imaging, or fusion.

\paragraph{Cross-modal transfer.}
The closest published setting is \emph{modality hallucination}~\cite{hoffman2016modalityhallucination}: training on one modality so the model can be queried at test time on another. RGB-to-thermal style translation~\cite{berg2018rgb2thermal,kniaz2018thermalgan} addresses the inverse direction, and domain-adversarial training~\cite{ganin2016domain} provides the general feature-alignment toolkit. None match the setup here, in which no cross-modal information is provided at training time and transfer emerges from shared low-level statistics (single-channel intensity, sky-background silhouettes). The result is presented as emergent in \S\ref{sec:grayscale}.

\paragraph{Cascaded rejection.}
The architectural pattern of an initial high-recall stage followed by progressively more discriminative rejection predates this application by two decades: \citet{viola2001rapid} for face detection, Cascade R-CNN~\cite{cai2018cascade} for general objects. Both reuse the same backbone across stages and differ from this design in two ways: (i) the rejection stages here are heterogeneous (XGBoost on per-frame scene features, then a separate MLP on detection ROI features), making it a refined-decision pipeline, not a refined-bounding-box pipeline; and (ii) the cascade boundary is the alert-emit decision (\S\ref{sec:alert_gate}), not the per-detection score. The Viola-Jones intent --- early rejection cheap, late rejection expensive --- is preserved.

\section{Comparison to Prior Work}
\label{sec:comparison}

Direct numerical comparison to other systems is unusually difficult: papers vary in dataset split, inference resolution, scoring rule (IoU/IoP, dual vs trust-aware), and headline metric (mAP, F1, success rate), and the scoring-rule audit alone swings F1 by 28~pp. The comparison below is in two parts: a qualitative architectural map (Table~\ref{tab:related_systems}) and a numerical panel for studies reporting comparable metrics on the same public benchmarks (Table~\ref{tab:numerical_comparison}).

\subsection{Architectural Comparison}

\begin{table}[htbp]
    \centering
    \caption{Architectural map of representative drone-detection systems. ``Modality'' covers what the system consumes; ``Discrimination'' covers how confusers are rejected. The cascade in this thesis is, to our knowledge, the only one that combines learned modality-trust fusion with a downstream per-frame distilled confuser verifier.}
    \label{tab:related_systems}
    \begin{tabular}{p{4.0cm}lll}
        \toprule
        System / paper & Modality & Discrimination & Notes \\
        \midrule
        Svanstr\"om 2021~\cite{svanstrom2021real} & RGB + IR (parallel) & Per-modality YOLO only & Dataset paper; no learned fusion \\
        Anti-UAV challenge baselines~\cite{jiang2021antiuav} & RGB + IR (parallel) & Trackers + YOLO; no confuser stage & Drone-tracking task, not raw detection \\
        Coluccia 2021 drone-vs-bird~\cite{coluccia2021dronevsbird} & RGB only & Single-stage YOLO retrained on bird negatives & Hard-negative training, no cascade \\
        Cross-modal attention fusion (e.g., \cite{wagner2016multispectral}) & RGB + IR (intermediate) & End-to-end CNN, no explicit confuser stage & Tightly coupled, requires paired training data \\
        \textbf{This thesis} & RGB + IR (decision-level) & Trust classifier + per-frame distilled (\texttt{mlp\_v5}) confuser verifier & Learned per-frame modality choice; cascade designed for fail-open \\
        \bottomrule
    \end{tabular}
\end{table}

Prior work uses either a single detector or a tightly-coupled multi-modal one and treats confuser discrimination as a training-time problem. This thesis explicitly separates the recall stage (detector) from the precision stage (trust classifier + per-frame verifier) and makes the modality-trust decision per frame.

\subsection{Numerical Comparison}

\begin{table}[htbp]
    \centering
    \caption{Numerical comparison on public benchmarks. \textbf{All comparisons require the caveats below the table.} ``Scoring'' indicates the matching rule used in the cited paper where it can be determined; where it cannot, the cell is marked \texttt{ND} (not disclosed).}
    \label{tab:numerical_comparison}
    \begin{tabular}{lllccc}
        \toprule
        Dataset & System & Scoring & $P$ & $R$ & F1 \\
        \midrule
        \multirow{2}{*}{Svanstr\"om (drone, RGB)} & Svanstr\"om 2021 paper (as reported) & \texttt{ND} (IoU?) & ND & ND & $\sim$0.93 \\
                                                   & \textbf{This thesis (baseline RGB, S1)} & IoP@0.5 & 0.940 & 0.961 & 0.950 \\
        \midrule
        \multirow{2}{*}{Svanstr\"om (drone, IR)}  & Svanstr\"om 2021 paper (as reported) & \texttt{ND} & ND & ND & $\sim$0.92 \\
                                                   & \textbf{This thesis (IR \texttt{v3b}, \texttt{imgsz=640})} & IoP@0.5 & 0.950 & 0.973 & 0.961 \\
        \midrule
        \multirow{2}{*}{Anti-UAV (drone, RGB)}    & Anti-UAV baselines (range) & IoU & ND & ND & 0.90--0.97 \\
                                                   & \textbf{This thesis (baseline RGB)} & IoU@0.5 & 0.992 & 0.995 & 0.994 \\
        \midrule
        \multirow{2}{*}{Anti-UAV (drone, IR)}     & Anti-UAV baselines (range) & IoU & ND & ND & 0.90--0.96 \\
                                                   & \textbf{This thesis (IR \texttt{v3b})} & IoU@0.5 & 0.987 & 0.945 & 0.965 \\
        \midrule
        Confuser zoo (no GT, RGB)                  & \emph{No published baseline} & ND & ND & ND & ND \\
                                                   & \textbf{This thesis (cascade S3)} & fire-rate & ND & ND & 0.8\% fire \\
        \bottomrule
    \end{tabular}
\end{table}

Three caveats apply to every row of Table~\ref{tab:numerical_comparison}.

\emph{Scoring rule.} The cited Svanstr\"om and Anti-UAV baselines rarely disclose IoU threshold and almost never disclose whether multi-modal scoring is dual or trust-aware. The 28-pp swing in \S\ref{sec:scoring_audit} is the order-of-magnitude of the resulting comparison noise. IoP@0.5, used here for Svanstr\"om RGB, is more lenient than IoU@0.5; an IoU@0.5 re-score would produce lower numbers.

\emph{Dataset split.} The Svanstr\"om paper~\cite{svanstrom2021real} does not enforce a canonical train/test split. This thesis evaluates on the full Svanstr\"om corpus (no Svanstr\"om data in the RGB training set) but cannot guarantee comparable difficulty to other splits. Anti-UAV has a more canonical split but uses sequence-level tracking metrics not directly equivalent to per-frame F1.

\emph{Task.} Anti-UAV baselines are tracking systems. Their state-accuracy numbers (0.90--0.97 range across years) include temporal aggregation absent from this thesis's frame-level numbers. The Anti-UAV row is included for orientation only.
% [source: ledger=antiuav-saturated; eval=rgb_antiuav_baseline; run=eval/ablate.py; cache=eval/results/_ablation/2026-05-18T22-48-19/master.csv; config=antiuav_iou_1280_s5]

Two qualitative claims survive these caveats. The per-modality numbers are in the same range as published Svanstr\"om results, establishing that the detectors are competitive baselines rather than weak strawmen. No published system reports a separately-evaluated OOD confuser-zoo fire rate at the granularity of Table~\ref{tab:cum_confuser}, because no published system explicitly separates a confuser-rejection stage from drone detection. The Stage-3 suppression numbers ($52.1\% \to 10.3\%$ with sa32; $52.1\% \to 0.8\%$ with \texttt{fusion\_no\_fn\_v1.1}; classifier-only $20.5\%$ and $1.6\%$ respectively; \S\ref{sec:cumulative}) therefore cannot be directly benchmarked. The contribution is the architectural separation that makes such numbers measurable.
% [source: ledger=trust-classifier-conditional; eval=clfzoo_sa32,clfzoo_fnfn; run=eval/cumulative_halluc.py; cache=eval/results/_cumulative_halluc/; config=confuser_zoo_1280]

% ======================================================================
% CHAPTER 3 — METHODOLOGY
% ======================================================================
